
\titulum{3}{Dans l'octave de la fête du Corps du Christ}{Propre du Temps}{Temps après la Pentecôte}

\rubrica{Aux petites heures, on peut prendre le ton de Noël des hymnes.}

\titulum{3}{Deuxième dimanche après la Pentecôte, dans l'octave de la fête du Corps du Christ}{Propre du Temps}{Temps après la Pentecôte}

\titulumrubrica{Aux premières Vêpres}

\capitulum{1 Jn. 3: 13-14}{Caríssimi: Nolíte mirári, si odit vos mundus.\pscross{} Nos scimus quóniam transláti sumus de morte ad vitam,\psstar{} quóniam dilígimus fratres.}{Bien-aimés, ne soyez pas étonnés si le monde a de la haine contre vous. Nous, nous savons que nous sommes passés de la mort à la vie, parce que nous aimons nos frères.}

\versiculus{Cibávit illos ex ádipe fruménti, allelúia.}{Et de petra, melle saturávit eos, allelúia.}{Il les a nourris de la fleur du froment, alléluia.}{Et les a rassasiés de miel pris au rocher, alléluia.}

\gscore{H1F7VAM}{Mag}{Le jeune Samuel servait le Seigneur en présence d'Élie et gardait précieusement la parole du Seigneur.}

\oratio{Sancti nóminis tui, Dómine, timórem páriter et amórem fac nos habére perpétuum:\pscross{} quia nunquam tua gubernatióne destítuis,\psstar{} quos in soliditáte tuæ dilectiónis instítuis. Per Dóminum.}{Fais, Seigneur, que nous ayons toujours la crainte et l'amour de ton saint Nom, parce que tu ne cesses jamais de diriger ceux que tu établis dans la solidité de ton amour. Par Jésus Christ.}	

\paragraph{Aux Laudes}

\rubrica{Capitule \normaltext{Caríssimi}, ci-dessus.}

\versiculus{Panem cæli dedit eis, allelúia.}{Panem Angelórum manducávit homo, allelúia.}{Il leur a donné le pain du ciel, alléluia.}{L'homme a mangé le pain des anges, alléluia.}

\gscore{H2F1LAB}{Ben}{Un homme fit un grand festin et invita beaucoup de convives ; et il envoya son serviteur, à l'heure du repas, dire aux invités de venir car tout est prêt, alléluia.}

\gscore{H2F1VAM}{Mag}{Va vite sur les places et dans les rues de la ville et invite à entrer les pauvres et les faibles, les aveugles et les boiteux, afin que ma maison soit remplie, alléluia.}

\titulum{3}{Troisième dimanche après la Pentecôte, dans l'octave de la fête du Sacré-Cœur}{Propre du Temps}{Après le Sacré-Cœur}

Protéctor in te sperántium, Deus, sine quo nihil est válidum, nihil sanctum; multíplica super nos misericórdiam tuam; ut, te rectóre, te duce, sic transeámus per bona temporália, ut non amittámus ætérna.
